\begin{figure}[htbp]
\centering
\begin{tikzpicture}

% ========== Panel (a): Original Unit Square ==========
\begin{scope}[xshift=0cm]
    \node[font=\bfseries] at (1.5, 3.8) {(a) Unit Square};

    % Axes
    \draw[->, thick] (-0.5, 0) -- (3.5, 0) node[right] {$x$};
    \draw[->, thick] (0, -0.5) -- (0, 3.5) node[above] {$y$};

    % Grid
    \draw[help lines, UofSC50Black!40] (-0.5, -0.5) grid (3, 3);

    % Unit square (filled)
    \fill[UofSCAtlantic!20] (0, 0) -- (1, 0) -- (1, 1) -- (0, 1) -- cycle;
    \draw[very thick, UofSCAtlantic] (0, 0) -- (1, 0) -- (1, 1) -- (0, 1) -- cycle;

    % Basis vectors
    \draw[->, line width=2pt, UofSCGarnet] (0, 0) -- (1, 0) node[below, font=\bfseries] {$\vece_1$};
    \draw[->, line width=2pt, UofSCHorseshoe] (0, 0) -- (0, 1) node[left, font=\bfseries] {$\vece_2$};

    % Corner labels
    \node[below left, font=\small] at (0, 0) {$O$};
    \node[circle, fill=UofSCBlack, minimum size=4pt, inner sep=0pt] at (0, 0) {};
    \node[circle, fill=UofSCBlack, minimum size=4pt, inner sep=0pt] at (1, 0) {};
    \node[circle, fill=UofSCBlack, minimum size=4pt, inner sep=0pt] at (1, 1) {};
    \node[circle, fill=UofSCBlack, minimum size=4pt, inner sep=0pt] at (0, 1) {};

    % Matrix label
    \node[align=center, font=\footnotesize] at (1.5, -1.2) {$\mI = \begin{bmatrix} 1 & 0 \\ 0 & 1 \end{bmatrix}$};
\end{scope}

% ========== Panel (b): Shear Transformation ==========
\begin{scope}[xshift=5cm]
    \node[font=\bfseries] at (2, 3.8) {(b) Shear};

    % Axes
    \draw[->, thick] (-0.5, 0) -- (4, 0) node[right] {$x$};
    \draw[->, thick] (0, -0.5) -- (0, 3.5) node[above] {$y$};

    % Grid
    \draw[help lines, UofSC50Black!40] (-0.5, -0.5) grid (3.5, 3);

    % Original unit square (ghost)
    \draw[thick, UofSC50Black, dashed] (0, 0) -- (1, 0) -- (1, 1) -- (0, 1) -- cycle;

    % Sheared parallelogram: A = [1 0.8; 0 1]
    % (0,0) -> (0,0), (1,0) -> (1,0), (0,1) -> (0.8,1), (1,1) -> (1.8,1)
    \fill[UofSCRose!20] (0, 0) -- (1, 0) -- (1.8, 1) -- (0.8, 1) -- cycle;
    \draw[very thick, UofSCRose] (0, 0) -- (1, 0) -- (1.8, 1) -- (0.8, 1) -- cycle;

    % Transformed basis vectors
    \draw[->, line width=2pt, UofSCGarnet] (0, 0) -- (1, 0) node[below, font=\bfseries] {$\mA\vece_1$};
    \draw[->, line width=2pt, UofSCHorseshoe] (0, 0) -- (0.8, 1) node[above left, font=\bfseries] {$\mA\vece_2$};

    % Corners
    \node[circle, fill=UofSCBlack, minimum size=4pt, inner sep=0pt] at (0, 0) {};
    \node[circle, fill=UofSCBlack, minimum size=4pt, inner sep=0pt] at (1, 0) {};
    \node[circle, fill=UofSCBlack, minimum size=4pt, inner sep=0pt] at (1.8, 1) {};
    \node[circle, fill=UofSCBlack, minimum size=4pt, inner sep=0pt] at (0.8, 1) {};

    % Matrix label
    \node[align=center, font=\footnotesize] at (2, -1.2) {$\mA = \begin{bmatrix} 1 & 0.8 \\ 0 & 1 \end{bmatrix}$};
\end{scope}

% ========== Panel (c): Rotation ==========
\begin{scope}[xshift=10cm]
    \node[font=\bfseries] at (1.5, 3.8) {(c) Rotation ($45^{\circ}$)};

    % Axes
    \draw[->, thick] (-1.5, 0) -- (2.5, 0) node[right] {$x$};
    \draw[->, thick] (0, -0.5) -- (0, 3.5) node[above] {$y$};

    % Grid
    \draw[help lines, UofSC50Black!40] (-1.5, -0.5) grid (2, 3);

    % Original unit square (ghost)
    \draw[thick, UofSC50Black, dashed] (0, 0) -- (1, 0) -- (1, 1) -- (0, 1) -- cycle;

    % Rotated square: R(45°) = [cos45 -sin45; sin45 cos45] = [0.707 -0.707; 0.707 0.707]
    % (0,0) -> (0,0), (1,0) -> (0.707, 0.707), (0,1) -> (-0.707, 0.707), (1,1) -> (0, 1.414)
    \fill[UofSCCongaree!20] (0, 0) -- (0.707, 0.707) -- (0, 1.414) -- (-0.707, 0.707) -- cycle;
    \draw[very thick, UofSCCongaree] (0, 0) -- (0.707, 0.707) -- (0, 1.414) -- (-0.707, 0.707) -- cycle;

    % Transformed basis vectors
    \draw[->, line width=2pt, UofSCGarnet] (0, 0) -- (0.707, 0.707) node[right, font=\bfseries] {$\mR\vece_1$};
    \draw[->, line width=2pt, UofSCHorseshoe] (0, 0) -- (-0.707, 0.707) node[left, font=\bfseries] {$\mR\vece_2$};

    % Rotation arc
    \draw[->, thick, UofSCAtlantic] (0.5, 0) arc (0:45:0.5);
    \node[UofSCAtlantic, font=\small] at (0.7, 0.25) {$45^{\circ}$};

    % Corners
    \node[circle, fill=UofSCBlack, minimum size=4pt, inner sep=0pt] at (0, 0) {};
    \node[circle, fill=UofSCBlack, minimum size=4pt, inner sep=0pt] at (0.707, 0.707) {};
    \node[circle, fill=UofSCBlack, minimum size=4pt, inner sep=0pt] at (0, 1.414) {};
    \node[circle, fill=UofSCBlack, minimum size=4pt, inner sep=0pt] at (-0.707, 0.707) {};

    % Matrix label
    \node[align=center, font=\footnotesize] at (0.5, -1.2) {$\mR = \begin{bmatrix} \cos 45^{\circ} & -\sin 45^{\circ} \\ \sin 45^{\circ} & \cos 45^{\circ} \end{bmatrix}$};
\end{scope}

% ========== Panel (d): Scaling ==========
\begin{scope}[xshift=0cm, yshift=-6.5cm]
    \node[font=\bfseries] at (2.5, 3.8) {(d) Non-uniform Scaling};

    % Axes
    \draw[->, thick] (-0.5, 0) -- (5, 0) node[right] {$x$};
    \draw[->, thick] (0, -0.5) -- (0, 3.5) node[above] {$y$};

    % Grid
    \draw[help lines, UofSC50Black!40] (-0.5, -0.5) grid (4.5, 3);

    % Original unit square (ghost)
    \draw[thick, UofSC50Black, dashed] (0, 0) -- (1, 0) -- (1, 1) -- (0, 1) -- cycle;

    % Scaled rectangle: S = [2 0; 0 0.5]
    % (0,0) -> (0,0), (1,0) -> (2,0), (0,1) -> (0,0.5), (1,1) -> (2,0.5)
    \fill[UofSCHorseshoe!20] (0, 0) -- (2, 0) -- (2, 0.5) -- (0, 0.5) -- cycle;
    \draw[very thick, UofSCHorseshoe] (0, 0) -- (2, 0) -- (2, 0.5) -- (0, 0.5) -- cycle;

    % Transformed basis vectors
    \draw[->, line width=2pt, UofSCGarnet] (0, 0) -- (2, 0) node[below, font=\bfseries] {$\mS\vece_1$};
    \draw[->, line width=2pt, UofSCAtlantic] (0, 0) -- (0, 0.5) node[left, font=\bfseries] {$\mS\vece_2$};

    % Corners
    \node[circle, fill=UofSCBlack, minimum size=4pt, inner sep=0pt] at (0, 0) {};
    \node[circle, fill=UofSCBlack, minimum size=4pt, inner sep=0pt] at (2, 0) {};
    \node[circle, fill=UofSCBlack, minimum size=4pt, inner sep=0pt] at (2, 0.5) {};
    \node[circle, fill=UofSCBlack, minimum size=4pt, inner sep=0pt] at (0, 0.5) {};

    % Annotations
    \node[UofSCGarnet, font=\footnotesize] at (3.2, 0.15) {stretch by 2};
    \node[UofSCAtlantic, font=\footnotesize] at (1.8, 1.2) {shrink by $\frac{1}{2}$};

    % Matrix label
    \node[align=center, font=\footnotesize] at (2.5, -1.2) {$\mS = \begin{bmatrix} 2 & 0 \\ 0 & 0.5 \end{bmatrix}$};
\end{scope}

% ========== Panel (e): Combined Transformation ==========
\begin{scope}[xshift=7cm, yshift=-6.5cm]
    \node[font=\bfseries] at (3, 3.8) {(e) Combined: Rotation + Scaling};

    % Axes
    \draw[->, thick] (-1.5, 0) -- (5, 0) node[right] {$x$};
    \draw[->, thick] (0, -1.5) -- (0, 3.5) node[above] {$y$};

    % Grid
    \draw[help lines, UofSC50Black!40] (-1.5, -1.5) grid (4.5, 3);

    % Original unit square (ghost)
    \draw[thick, UofSC50Black, dashed] (0, 0) -- (1, 0) -- (1, 1) -- (0, 1) -- cycle;

    % Combined: Scale by [1.5, 0.8] then rotate 30°
    % First scale: (1,0) -> (1.5, 0), (0,1) -> (0, 0.8)
    % Then rotate 30°: [cos30 -sin30; sin30 cos30]
    % R*S*e1 = R*(1.5, 0) = (1.5*cos30, 1.5*sin30) = (1.299, 0.75)
    % R*S*e2 = R*(0, 0.8) = (-0.8*sin30, 0.8*cos30) = (-0.4, 0.693)
    % (1,1) -> R*S*(1,1) = R*(1.5, 0.8) = (1.5*cos30 - 0.8*sin30, 1.5*sin30 + 0.8*cos30) = (0.899, 1.443)
    \fill[UofSCGarnet!20] (0, 0) -- (1.299, 0.75) -- (0.899, 1.443) -- (-0.4, 0.693) -- cycle;
    \draw[very thick, UofSCGarnet] (0, 0) -- (1.299, 0.75) -- (0.899, 1.443) -- (-0.4, 0.693) -- cycle;

    % Transformed basis vectors
    \draw[->, line width=2pt, UofSCAtlantic] (0, 0) -- (1.299, 0.75) node[right, font=\bfseries] {$\mM\vece_1$};
    \draw[->, line width=2pt, UofSCHorseshoe] (0, 0) -- (-0.4, 0.693) node[above left, font=\bfseries] {$\mM\vece_2$};

    % Corners
    \node[circle, fill=UofSCBlack, minimum size=4pt, inner sep=0pt] at (0, 0) {};
    \node[circle, fill=UofSCBlack, minimum size=4pt, inner sep=0pt] at (1.299, 0.75) {};
    \node[circle, fill=UofSCBlack, minimum size=4pt, inner sep=0pt] at (0.899, 1.443) {};
    \node[circle, fill=UofSCBlack, minimum size=4pt, inner sep=0pt] at (-0.4, 0.693) {};

    % Matrix label
    \node[align=center, font=\footnotesize] at (3, -1.2) {$\mM = \mR(30^{\circ}) \cdot \mS = \begin{bmatrix} 1.30 & -0.40 \\ 0.75 & 0.69 \end{bmatrix}$};
\end{scope}

\end{tikzpicture}
\caption{Geometric interpretation of $2 \times 2$ matrix transformations acting on the unit square. (a) The identity matrix preserves all vectors. (b) Shear transforms the square into a parallelogram by sliding points parallel to one axis. (c) Rotation preserves lengths and angles (orthogonal transformation). (d) Non-uniform scaling stretches or compresses along each axis independently. (e) Combined transformations compose: here scaling followed by rotation. In each case, the matrix columns represent the transformed basis vectors $\mA\vece_1$ and $\mA\vece_2$. The dashed squares show the original unit square for reference.}
\label{fig:matrix_transformation}
\end{figure}
